% =====================================================================
% Responsabilidades de cada integrante --- Brunno
% Só o Brunno edita este arquivo. Conferir na revisão final (o que cada um fez de fato).
% =====================================================================
\section{Responsabilidades}\label{sec:responsabilidades}

A Tabela~\ref{tab:responsabilidades} mostra o que cada integrante fez no código, no relatório e
nos documentos da Etapa 1. Os commits de cada um estão no histórico do repositório Git do grupo.
Todos participaram da gravação do vídeo e da apresentação.

\begin{table}[htbp]
\centering
\caption{Responsabilidades de cada integrante.}
\label{tab:responsabilidades}
\small
\begin{tabularx}{\textwidth}{@{}l Y Y@{}}
\toprule
Integrante & Código & Relatório e Etapa 1 \\
\midrule
Brunno & Classe \codigo{Grafo} e reindexação; leitores das quatro bases; geradores Erdős--Rényi,
  Barabási--Albert e Watts--Strogatz; organização do repositório e revisão dos PRs
  & Introdução, bases, modelagem e responsabilidades; montagem e revisão do relatório. Etapa 1:
  modelagem, diagramas de classes e de pacotes e montagem do documento \\
Vitória & Grau e densidade; BFS, DFS, componentes e Dijkstra; distância média, diâmetro,
  proximidade e harmônica
  & Trabalhos relacionados; grau, caminhos e proximidade; roteiro do vídeo. Etapa 1: plano de
  desenvolvimento \\
Iago & Intermediação de Brandes (exata e amostrada), PageRank e número de núcleo; medições de
  desempenho e comparação com o NetworkX
  & Centralidades e avaliação de desempenho; edição do vídeo. Etapa 1: diagrama de sequência \\
Arthur & Triângulos, agrupamento e transitividade; modularidade, Louvain e propagação de rótulos;
  NMI e ARI
  & Agrupamento e comunidades; limitações e ética; slides. Etapa 1: diagrama de atividades do
  Louvain \\
Ravi & Limiar epidêmico, simulação SIR, poder de espalhamento, $\tau$ de Kendall e comparação dos
  rankings
  & Problema aplicado e resultados dos influenciadores. Etapa 1: descrição do problema \\
Kaio & Posicionamento de Fruchterman--Reingold e interface web
  & Interface e conclusão. Etapa 1: casos de uso e esboço das telas \\
\bottomrule
\end{tabularx}
\end{table}
